\medskip

\begin{tabularx}{\linewidth}{|X|}\hline
\textbf{Principaux domaines abordés :} Suites;
Fonctions, Fonction exponentielle.\\ \hline
\end{tabularx}

\bigskip

\textbf{Partie A}

\medskip

On considère la fonction $f$ définie pour tout réel $x$ par:

\[f(x) = 1+x - \text{e}^{0,5x - 2}.\]

On admet que la fonction $f$ est dérivable sur $\R$. On note $f'$ sa dérivée.

\medskip

\begin{enumerate}
\item 
	\begin{enumerate}
		\item 
On a $\displaystyle\lim_{x \to - \infty} 0,5x = - \infty$, donc $\displaystyle\lim_{x \to - \infty} 0,5x - 2 = - \infty$ et par composition de limites avec $X = 0,5x - 2$, $\displaystyle\lim_{X \to - \infty} \text{e}^X = 0$ et comme $\displaystyle\lim_{x \to - \infty} x = - \infty$ par somme de limites :

$\displaystyle\lim_{x \to - \infty} f(x) = - \infty$.
		\item Démontrer que, pour tout réel $x$ non nul, $f(x) = 1+ 0,5x\left(2 - \dfrac{\text{e}^{0,5x}}{0,5x} \times \text{e}^{-2}\right)$.

On a $f(x) = 1 + \left[x - \text{e}^{0,5x - 2}\right]$ et en factorisant $0,5x$ dans le crochet :

$f(x)  = 1 + 0,5x\left[2 - \dfrac{\text{e}^{0,5x - 2}}{0,5x} \right] =  1+ 0,5x\left(2 - \dfrac{\text{e}^{0,5x}}{0,5x} \times \text{e}^{-2}\right)$.

En déduire la limite de la fonction $f$ en $+\infty$.

Avec $X = 0,5x$, on a $\displaystyle\lim_{x \to - \infty}0,5x = + \infty$, d'où $\displaystyle\lim_{x \to - \infty} X = + \infty$.

mais $\displaystyle\lim_{x \to + \infty}\dfrac{\text{e}^X}{X} = + \infty$ et $\displaystyle\lim_{x \to + \infty}\dfrac{\text{e}^X}{X} \times \text{e}^{-2} = + \infty$, puis $\displaystyle\lim_{x \to + \infty}-\dfrac{\text{e}^X}{X} \times \text{e}^{-2} = - \infty$ et enfin

$\displaystyle\lim_{x \to + \infty}2 -\dfrac{\text{e}^X}{X} \times \text{e}^{-2} = - \infty$

Enfin par produit de limites :

$\left\{\begin{array}{l}
\displaystyle\lim_{x \to + \infty} 0,5x = + \infty\\
\displaystyle\lim_{x \to + \infty}2 -\dfrac{\text{e}^X}{X} \times \text{e}^{-2} = - \infty
\end{array}\right. \Rightarrow \displaystyle\lim_{x \to + \infty}f(x) = - \infty$.

	\end{enumerate}
\item 
	\begin{enumerate}
		\item 
$f$ est une somme de fonctions dérivables sur $\R$ et sur cet intervalle :

$f'(x) = 1 - 0,5\text{e}^{0,5x - 2}$.
		\item 
On a dans $\R$,

$f'(x) < 0 \iff 1 - 0,5\text{e}^{0,5x - 2} < 0 \iff 1 < 0,5\text{e}^{0,5x - 2} \iff 2 < \text{e}^{0,5x - 2}$. par croissance sur $\R_{+}^*$ de la fonction logarithme népérien $ \ln 2 < 0,5x - 2 \iff 2 + \ln 2 < 0,5x \iff 4 + 2\ln 2 < x$.

Conclusion : $S = ]4 + 2\ln (2)~;~+\infty[$

	\end{enumerate}
\item 
On démontrerait de la même façon que $f'(x) > 0$ sur l'intervalle $]- \infty~;~4 + 2\ln 2[$.

La fonction $f$ est donc croissante sur $]- \infty~;~4 + 2\ln 2[$ puis décroissante sur $]4 + 2\ln (2)~;~+\infty[$.

$f\left(4 + 2\ln 2 \right) = 1 + 4 + 2\ln 2 - \text{e}^{2 + \ln 2 - 2} = 5 + 2\ln 2 - \text{e}^{\ln 2} = 5 + 2\ln 2 - 2 = 3 + 2\ln 2$. C'est le maximum de la fonction $f$ sur $\R$.

Comme $\displaystyle\lim_{x \to - \infty} x = - \infty$ et $\displaystyle\lim_{x \to - \infty} \text{e}^{0,5x - 2} = 0$, on a par somme de limites  $\displaystyle\lim_{x \to - \infty} f(x) = - \infty$.
\item 
L'intervalle $[-1~;~0]$ est inclus dans l'intervalle $]- \infty~;~4 + 2\ln 2[$, donc est croissante sur cet intervalle de $f(-1) = 1 - 1 - \text{e}^{-1,5} = - \text{e}^{-1,5} < 0$ à $f(0) = 1 - \text{e}^{-2} \approx 0,85 > 0$.

La fonction $f$ étant continue puisque dérivable le théorème des valeurs intermédiaires montre qu'il existe un réel unique $\alpha \in[-1~;~0]$ tel que $f(\alpha) = 0$.
\end{enumerate}

\bigskip

\textbf{Partie B}

\medskip

On considère la suite $\left(u_n\right)$ définie par $u_0 = 0$ et, pour tout entier naturel $n$,
\: $u_{n+1} = f\left(u_n\right)$ où $f$ est la fonction définie à la partie A.

\medskip

\begin{enumerate}
\item 
	\begin{enumerate}
		\item Démontrer par récurrence que, pour tout entier naturel $n$, on a : 

		\[u_n \leqslant u_{n+1} \leqslant 4.\]
\emph{Initialisation } : On a $u_0 = 0$, puis $u_1 = f\left(u_0\right) = f(0) = 1 - \text{e}^{-2} \approx 0,85$.

On a donc bien pour $n= 0$ : $u_0 \leqslant u_1 \leqslant 4$.

\emph{Hérédité} : soit $n \in \N$ et supposons que $u_n \leqslant u_{n+1} \leqslant 4$.

On a vu dans la partie A que sur l'intervalle $]- \infty~;~4 = 2\ln 2[$, donc sur l'intervalle $]- \infty~;~4[$, la fonction est croissante et l'hypothèse de récurrence donne

$f\left(u_n \right) \leqslant f\left(u_{n+1} \right) \leqslant f(4)$.

Or $f(4) = 1 + 4 - \text{e}^{2 - 2} = 5 -1 = 4$.

On a donc $u_{n+1} \leqslant u_{n+2} \leqslant 4$ : l'encadrement est vrai au rang $n + 1$.

Conclusion : l'encadrement est vrai au rang 0, et s'il est vrai pour $n \in \N$, il l'est aussi pour $n + 1$.
Le principe de récurrence montre donc que pour tout $n \in \N$, \: $u_n \leqslant u_{n+1} \leqslant 4$.
		
		\item 

Le résultat précédent montre que la suite $\left(u_n\right)$ est croissante et majorée par 4 : elle donc convergente vers une limite $\ell$ telle que $\ell \leqslant 4$.
	\end{enumerate}	
\item 
	\begin{enumerate}
		\item 
		

$\ell = f(\ell) \iff \ell = 1 + \ell - \text{e}^{0,5\ell - 2} \iff \text{e}^{0,5\ell - 2} = 1 \iff 0,5\ell - 2 = 0$ (par croissance de la fonction logarithme népérien sur $\R_{+}^*$), d'où $0,5\ell = 2 \iff \ell = 4$.

La suite $\left(u_n\right)$ a pour limite 4.
		\item ~
		

















L'algorithme calcule les premiers termes de la suite $\left(u_n\right)$ jusqu'à celui qui est supérieur à 3,99 : 12 signifie que $u_{12}$ est le premier terme supérieur à 3,99.
	\end{enumerate}
\end{enumerate}

\bigskip
